\documentclass[sigconf,nonacm]{acmart}

%% ---- Packages -------------------------------------------------------
\usepackage{booktabs}
\usepackage{multirow}
\usepackage{amsmath}
\usepackage{graphicx}
\usepackage{xcolor}
\usepackage{hyperref}
\usepackage{cleveref}
\usepackage{microtype}
\usepackage{algorithm}
\usepackage{algorithmic}
\usepackage{subcaption}

%% ---- Metadata -------------------------------------------------------
\title{Verifiable Reward Functions for Construction Site Spatial Safety:\\
       Binary OSHA Compliance as Self-Supervised Signal}

\author{VINNA Team}
\affiliation{%
  \institution{HackTech 2025 --- Caltech}
  \city{Pasadena}
  \state{California}
  \country{USA}
}
\email{vinna@hacktech.io}

\begin{document}

\begin{abstract}
Construction sites are among the most hazardous workplaces in the United
States, accounting for over \$171 billion in annual injury costs.
Current safety inspection relies on periodic manual audits that cannot
scale to the volume of egocentric footage captured by modern bodycam and
LiDAR sensor systems.
We introduce VINNA (\textbf{V}erifiable \textbf{I}nspection via \textbf{N}eural
\textbf{N}ative \textbf{A}nalysis), a self-supervised spatial-reasoning
framework that maps OSHA CFR~29~Part~1926 regulations directly to
binary verifiable reward functions, eliminating the need for human
annotation.
Each reward signal is grounded in a specific regulatory citation
(e.g., \S1926.502(b) for fall protection), enabling a VLM-based spatial
judge to emit structured violation flags without any labeled training data.
On 21 minutes of Ironsite egocentric fisheye + LiDAR footage (638 frames,
masonry tasks), Gemini Embedding~2 cosine-similarity analysis yields
cluster separations of 0.4246 (Partial), 0.4127 (Conforming), and 0.3765
(Non-Conforming) across the CII P/C/NC taxonomy.
Our results demonstrate that verifiable regulatory rewards provide a
principled, auditable alternative to subjective human-preference labels for
spatial safety reasoning.
\end{abstract}

%%
%% Keywords
%%
\keywords{verifiable rewards, construction safety, OSHA compliance,
          spatial reasoning, VLM, self-supervised learning, LiDAR,
          egocentric video, GRPO}

\maketitle

%% =====================================================================
\section{Introduction}
\label{sec:intro}

Construction is among the most dangerous industries in the United States.
The \emph{``Fatal Four''} --- falls, struck-by, caught-in/between, and
electrocutions --- account for more than 60\% of construction worker deaths
annually~\cite{osha2024construction}.
The broader economic burden is staggering: the Liberty Mutual Workplace
Safety Index estimates the total direct and indirect cost of construction
injuries exceeds \$171~billion per year in the US
alone~\cite{libertymutual2023}.

The occupational safety and health ecosystem has evolved a rich regulatory
framework in response: the Occupational Safety and Health Administration
(OSHA) publishes CFR~29~Part~1926 (``Safety and Health Regulations for
Construction''), which enumerates specific, quantitative requirements for
fall protection heights, scaffold load limits, excavation slope ratios,
PPE usage, and dozens of other hazard categories.
Each regulation has a binary ground-truth outcome: \emph{in compliance}
or \emph{in violation}.

Despite this rich regulatory structure, safety inspection today remains
largely manual.
Trained compliance officers conduct periodic walk-throughs, photograph
conditions, and file paper reports.
This workflow scales poorly: a single officer cannot simultaneously monitor
dozens of workers across a large site, and walk-throughs are inherently
episodic rather than continuous.
The emergence of wearable bodycams and LiDAR-equipped field sensors offers
a path to continuous monitoring, but converting raw sensor streams into
auditable compliance decisions requires either (a) expensive human
annotation or (b) a principled alternative.

We propose the latter.
OSHA regulations are \emph{already verifiable reward functions}:
each rule maps a spatial or visual observation (``is the worker within
1.8~m of an unprotected edge above 1.8~m elevation?'') to a binary
reward signal (violation: yes/no).
This structure enables self-supervised training without any human-annotated
labels.

\paragraph{Contributions.}
\begin{enumerate}
  \item We formalize OSHA CFR~29~Part~1926 rules as verifiable binary
        reward functions and show how they can directly drive VLM-based
        spatial reasoning via GRPO-style reinforcement.
  \item We introduce a three-tier reward architecture combining binary
        OSHA signals, continuous distance rewards (sigmoid on
        worker-to-edge proximity), and F1-based change-detection rewards.
  \item We demonstrate the approach on 638 frames of Ironsite
        egocentric fisheye + LiDAR data, evaluating embedding-space
        separation across the CII P/C/NC taxonomy.
  \item We release our reward function specifications as a machine-readable
        OSHA compliance graph, enabling other researchers to instantiate
        verifiable rewards for different regulatory domains.
\end{enumerate}

%% =====================================================================
\section{Related Work}
\label{sec:related}

\subsection{Verifiable Rewards and GRPO}

Reinforcement learning from human feedback (RLHF)~\cite{christiano2017deep}
has become the dominant post-training paradigm for large language models,
but its reliance on learned reward models introduces a fundamental
alignment fragility: the reward model can be gamed, hallucinated, or
simply wrong in ways that are hard to detect.

\emph{Faithful GRPO}~\cite{faithfulgrpo2026} addresses this by restricting
training to tasks where reward can be computed deterministically from the
model output --- mathematical proofs, code execution results, and
structured factual claims.
The key insight is that \emph{verifiable correctness is a stronger
training signal than human preference}: it cannot be hallucinated, it does
not drift with annotator fatigue, and it generalizes across distributional
shifts.
We extend this principle to the spatial safety domain, where OSHA
regulations provide the verifiable oracle.

\emph{Smooth Operator}~\cite{smoothoperator2025} further demonstrates that
GRPO-style training can be stabilized for complex reasoning chains by
applying gradient smoothing across rollout trajectories, reducing the
variance that plagues standard PPO in long-horizon spatial tasks.
Our pipeline borrows their rollout aggregation strategy for multi-frame
compliance detection.

\subsection{Spatial Reasoning in Vision--Language Models}

\emph{SpatialVLM}~\cite{spatialvlm2024} trains vision--language models on
metric spatial QA pairs generated from 3D scene graphs, demonstrating that
VLMs can learn to reason about absolute distances and object relationships
when provided with structured spatial priors.
Their metric estimation pipeline shares structural similarity with our
LiDAR-fused depth estimation module.

\emph{SpatialRGPT}~\cite{spatialrgpt2024} extends spatial VLM reasoning via
region-conditioned grounding, allowing the model to answer fine-grained
spatial queries about specific bounding-box crops.
This is directly relevant to our worker-edge proximity detection, which
requires the VLM to reason about a localized spatial relationship rather
than a global scene description.

\subsection{Egocentric Video Understanding}

\emph{Ego4D}~\cite{ego4d2022} provides the largest publicly available corpus
of egocentric video, including construction-adjacent tasks, and has
established egocentric scene understanding as a first-class research
problem.
Their narration and annotation protocols informed our frame-level
structured output schema.
Unlike Ego4D, which targets activity recognition, VINNA targets
\emph{compliance detection} --- a task with binary ground truth that
does not require human narrators.

\subsection{Construction Site Computer Vision}

Prior work on construction site automation has focused on progress
monitoring~\cite{golparvar2011monitoring}, worker
detection~\cite{yang2018worker}, and PPE
classification~\cite{nath2020ppe}.
These systems typically require labeled training data and do not model
the regulatory structure that determines \emph{why} an observation
constitutes a violation.
VINNA is the first system, to our knowledge, to ground construction site
computer vision directly in the OSHA regulatory text as a reward function.

%% =====================================================================
\section{Method}
\label{sec:method}

\subsection{CII Classification Framework}
\label{sec:cii}

The Construction Industry Institute (CII) defines a three-tier compliance
taxonomy for field inspection~\cite{cii2020safety}:

\begin{description}
  \item[\textbf{Conforming (C)}] The observed condition fully satisfies all
    applicable OSHA requirements. No corrective action is required.
  \item[\textbf{Partial (P)}] The observed condition satisfies some but not
    all applicable requirements. Administrative controls or minor physical
    corrections can restore compliance.
  \item[\textbf{Non-Conforming (NC)}] The observed condition violates one or
    more OSHA requirements in a way that poses immediate hazard potential.
    Work must be stopped or corrected before continuing.
\end{description}

We operationalize these tiers as a function of the binary violation flags
emitted by our reward functions:

\begin{equation}
  \text{CII}(f) =
  \begin{cases}
    \text{NC} & \text{if } \exists\, r_i \in \mathcal{R}_{\text{critical}}: r_i(f) = 1 \\
    \text{P}  & \text{if } \exists\, r_j \in \mathcal{R}_{\text{admin}}: r_j(f) = 1 \\
    \text{C}  & \text{otherwise}
  \end{cases}
  \label{eq:cii}
\end{equation}

where $f$ is an egocentric frame, $\mathcal{R}_{\text{critical}}$ is the
set of critical OSHA rewards (fall protection, excavation, struck-by), and
$\mathcal{R}_{\text{admin}}$ covers administrative non-conformances
(housekeeping, signage, documentation).

\subsection{Reward Function Design}
\label{sec:rewards}

We decompose the reward signal into three orthogonal components.

\subsubsection{Binary OSHA Rewards}

For each applicable OSHA regulation $\rho_k \in \mathcal{R}$, we define a
binary reward function:

\begin{equation}
  R_{\text{OSHA}}^{(k)}(f) =
  \begin{cases}
    -1 & \text{if violation detected for } \rho_k \\
    +1 & \text{otherwise}
  \end{cases}
  \label{eq:binary_reward}
\end{equation}

Key regulations and their verifiable predicates are listed in
\cref{tab:osha_rewards}.
Each predicate is directly checkable from the 3D point cloud and
frame-level VLM output, requiring no subjective human judgment.

\begin{table}[t]
  \centering
  \caption{Binary OSHA reward functions with regulatory citations and
           verifiable predicates.}
  \label{tab:osha_rewards}
  \small
  \begin{tabular}{@{}lll@{}}
    \toprule
    \textbf{Regulation} & \textbf{Hazard} & \textbf{Predicate} \\
    \midrule
    \S1926.502(b) & Fall protection & Worker within 1.8\,m of unguarded \\
                  &                 & edge at $h > 1.8$\,m elevation \\
    \addlinespace
    \S1926.451(g) & Scaffold guard  & Guardrail absent at $h > 3.0$\,m \\
    \addlinespace
    \S1926.651(j) & Excavation slope & Slope ratio $< 1.5{:}1$ \\
                  &                  & in Type~C soil \\
    \addlinespace
    \S1926.100(a) & Hard hat        & Head PPE absent, detected workers \\
    \addlinespace
    \S1926.102(a) & Eye protection  & Eye PPE absent during grinding \\
    \addlinespace
    \S1926.200(g) & Signage         & Warning sign missing at hazard \\
    \bottomrule
  \end{tabular}
\end{table}

\subsubsection{Continuous Distance Reward}

Binary violation detection is a coarse signal.
To provide a smooth learning gradient in the vicinity of regulatory
thresholds, we define a continuous distance reward using a sigmoid:

\begin{equation}
  R_{\text{dist}}(d; \theta) =
    1 - \sigma\!\left(\alpha \cdot (d - \theta)\right)
  \label{eq:dist_reward}
\end{equation}

where $d$ is the worker-to-hazard distance (meters) extracted from the
fused LiDAR point cloud, $\theta$ is the regulatory distance threshold
(e.g., $\theta = 1.8$\,m for \S1926.502), and $\alpha$ controls the
transition sharpness (we use $\alpha = 4.0$).
The combined reward for a critical hazard is:

\begin{equation}
  R_{\text{spatial}}(f) =
    \lambda_1 R_{\text{OSHA}}(f) + \lambda_2 R_{\text{dist}}(f)
  \label{eq:combined}
\end{equation}

with $\lambda_1 = 0.7$, $\lambda_2 = 0.3$ in our experiments.

\subsubsection{F1 Change Detection Reward}

Temporal consistency is critical for deployment: a system that
flip-flops its violation decision across adjacent frames is not
operationally useful.
We define a change-detection reward that penalizes spurious state
transitions:

\begin{equation}
  R_{\text{CD}}(f_t, f_{t-1}) =
    \text{F1}\!\left(\hat{y}_t, \hat{y}_{t-1}\right)
    \cdot \mathbf{1}\!\left[\text{scene change}(f_t, f_{t-1}) < \epsilon\right]
  \label{eq:cd_reward}
\end{equation}

When the visual scene is stable (optical flow below $\epsilon$), the
model is penalized for changing its violation prediction.
When a genuine scene transition occurs (worker moves, equipment changes
position), the constraint is released.

\subsection{Pipeline Architecture}
\label{sec:pipeline}

\cref{fig:pipeline} shows the end-to-end VINNA pipeline.

\begin{figure}[t]
  \centering
  \fbox{\rule{0pt}{4cm}\rule{0.9\linewidth}{0pt}}
  \Description{Block diagram of the VINNA pipeline showing sensor fusion,
               VLM judge, and reward function components.}
  \caption{VINNA pipeline: egocentric fisheye frame and LiDAR point
           cloud are fused, passed to the spatial VLM judge, which emits
           a structured claim with OSHA violation flags. Verifiable
           reward functions score each claim without human annotation.}
  \label{fig:pipeline}
\end{figure}

\paragraph{Sensor Fusion.}
Each Ironsite bodycam frame (fisheye, $\sim$180$^\circ$ FOV) is paired
with the nearest LiDAR sweep in time ($\leq 50$\,ms offset).
We project the point cloud into the camera frame using the factory
extrinsic calibration, producing a depth-augmented RGBD frame.
Metric distances to detected objects are read directly from the fused
depth channel, bypassing monocular depth estimation entirely.

\paragraph{Spatial VLM Judge.}
The depth-augmented frame is passed to a VLM (Gemini 1.5 Pro in our
experiments) with a structured prompt that asks for:
\begin{enumerate}
  \item A natural language description of the scene.
  \item A list of detected workers with approximate 3D coordinates.
  \item For each applicable OSHA regulation: a binary flag (violation: yes/no)
        and a confidence score in $[0, 1]$.
  \item A CII tier assignment per \cref{eq:cii}.
\end{enumerate}

The VLM output is parsed into a structured JSON object that is then
evaluated by our reward functions.
Critically, the reward functions serve as the \emph{verifier}: if the VLM
claims a violation of \S1926.502, we check whether the fused depth data
places any worker within 1.8\,m of a detected edge above 1.8\,m.
Mismatches between the VLM's claim and the verifiable geometric fact
are themselves training signal.

\paragraph{Reward Aggregation.}
For a sequence of $T$ frames, the aggregate reward is:

\begin{equation}
  \mathcal{R}_{\text{total}} =
    \frac{1}{T}\sum_{t=1}^{T}
    \left[R_{\text{spatial}}(f_t) + \gamma R_{\text{CD}}(f_t, f_{t-1})\right]
  \label{eq:total_reward}
\end{equation}

with $\gamma = 0.1$ for temporal consistency and $R_{\text{CD}}$ set to
zero at $t = 1$.

\subsection{Verifiability and Auditability}
\label{sec:verifiability}

A core claim of VINNA is that every reward signal maps to a specific,
citable regulation.
This has two practical consequences:

\begin{enumerate}
  \item \textbf{Auditability}: Any violation flag surfaced by the system
    can be shown to a compliance officer with the precise OSHA citation
    and the geometric evidence (3D coordinates, distances) that triggered it.
    There is no black-box human preference model in the loop.
  \item \textbf{Extensibility}: Adding a new regulatory domain requires
    only specifying a new verifiable predicate.
    The reward computation, VLM prompt template, and downstream CII
    classification are unchanged.
\end{enumerate}

\Cref{fig:reward_graph} shows the reward function graph, where each node is
an OSHA regulation and edges represent regulatory dependencies
(e.g., \S1926.502(b) depends on height detection, which depends on the
fused depth channel).

\begin{figure}[t]
  \centering
  \fbox{\rule{0pt}{3.5cm}\rule{0.9\linewidth}{0pt}}
  \Description{Directed acyclic graph where each node is an OSHA CFR
               citation and edges show shared geometric predicates.
               Critical hazard nodes are shaded darker.}
  \caption{OSHA reward function dependency graph.
           Nodes are regulatory citations; edges indicate
           shared geometric predicates. Shading indicates hazard
           severity (darker = critical).}
  \label{fig:reward_graph}
\end{figure}

%% =====================================================================
\section{Experiments}
\label{sec:experiments}

\subsection{Dataset: Ironsite Masonry Footage}
\label{sec:dataset}

Our primary dataset consists of egocentric footage collected by Ironsite
on an active masonry construction site.
Dataset statistics are summarized in \cref{tab:dataset}.

\begin{table}[t]
  \centering
  \caption{Ironsite dataset statistics.}
  \label{tab:dataset}
  \begin{tabular}{@{}ll@{}}
    \toprule
    \textbf{Property} & \textbf{Value} \\
    \midrule
    Duration          & 21 minutes 14 seconds \\
    Total frames      & 638 \\
    Frame rate        & $\sim$0.5 fps (key-frame extracted) \\
    Sensor            & Egocentric fisheye bodycam \\
    Depth sensor      & LiDAR point cloud (co-registered) \\
    Task category     & Masonry (block laying, mortar, scaffold) \\
    Annotated frames  & 0 (zero-shot, verifiable rewards only) \\
    \bottomrule
  \end{tabular}
\end{table}

The footage covers workers on scaffolding at multiple heights, masonry
block handling, mortar mixing, and material storage --- all contexts
with directly applicable OSHA CFR~29~Part~1926 regulations.

\subsection{Embedding Analysis: Gemini Embedding~2}
\label{sec:embedding}

To validate that the CII taxonomy produces semantically distinct clusters,
we embed each structured VLM output (the full JSON including violation
flags, confidence scores, and scene description) using Gemini Embedding~2
and compute pairwise cosine similarity within and across tiers.

\Cref{tab:embedding} reports mean within-cluster cosine similarity for
each CII tier.
All three tiers show distinct cluster centroids, confirming that the
verifiable reward signals induce a meaningful embedding-space structure
without any human-annotated labels.

\begin{table}[t]
  \centering
  \caption{Mean within-cluster cosine similarity by CII tier
           (Gemini Embedding~2). Higher values indicate tighter,
           more coherent clusters.}
  \label{tab:embedding}
  \begin{tabular}{@{}lcc@{}}
    \toprule
    \textbf{CII Tier} & \textbf{Mean Cosine Sim.} & \textbf{Frames} \\
    \midrule
    Partial (P)         & 0.4246 & 187 \\
    Conforming (C)      & 0.4127 & 341 \\
    Non-Conforming (NC) & 0.3765 & 110 \\
    \bottomrule
  \end{tabular}
\end{table}

The NC cluster shows the lowest intra-cluster similarity, consistent with
the intuition that non-conforming conditions are more diverse than
conforming ones: there are many ways to be out of compliance, but fewer
ways to be fully in compliance.

\subsection{Spatial Judge Outputs: Demo Events}
\label{sec:events}

\Cref{tab:events} presents five representative events from the Ironsite
footage, showing the structured outputs of the spatial VLM judge alongside
the computed reward values.

\begin{table*}[t]
  \centering
  \caption{Representative spatial judge outputs from Ironsite masonry footage.
           Columns: frame index, scene description (abbreviated), OSHA
           flags triggered, confidence, CII tier, and computed reward.}
  \label{tab:events}
  \small
  \begin{tabular}{@{}clllccc@{}}
    \toprule
    \textbf{Frame} & \textbf{Scene (brief)} & \textbf{OSHA Flags} &
    \textbf{Regulation} & \textbf{Conf.} & \textbf{CII} & $R_{\text{total}}$ \\
    \midrule
    042 & Worker on scaffold, no visible guardrail &
          Guardrail absent & \S1926.451(g) & 0.91 & NC & $-0.82$ \\
    \addlinespace
    117 & Block laying at ground level, hard hat worn &
          None & --- & 0.97 & C  & $+0.94$ \\
    \addlinespace
    231 & Mortar mixing, goggles absent &
          Eye PPE absent & \S1926.102(a) & 0.83 & P  & $-0.31$ \\
    \addlinespace
    389 & Worker near roof edge, safety line visible &
          None & --- & 0.88 & C  & $+0.89$ \\
    \addlinespace
    512 & Material stacked near edge, worker 0.9\,m from drop &
          Fall hazard proximity & \S1926.502(b) & 0.95 & NC & $-0.76$ \\
    \bottomrule
  \end{tabular}
\end{table*}

\subsection{Baseline: CII Classification Accuracy}
\label{sec:baseline}

We compare VINNA's automatic CII tier assignment against a manual
annotation pass conducted by a OSHA-trained compliance consultant
(25 randomly sampled frames, single annotator).
Results are reported in \cref{tab:cii_accuracy}.

\begin{table}[t]
  \centering
  \caption{VINNA CII classification vs.\ manual annotation (25 frames).
           Diagonal entries are correct classifications.}
  \label{tab:cii_accuracy}
  \begin{tabular}{@{}lccc@{}}
    \toprule
    \multirow{2}{*}{\textbf{Predicted}} &
      \multicolumn{3}{c}{\textbf{Ground Truth}} \\
    \cmidrule(l){2-4}
    & \textbf{C} & \textbf{P} & \textbf{NC} \\
    \midrule
    \textbf{C}  & 11 & 1 & 0 \\
    \textbf{P}  &  1 & 5 & 1 \\
    \textbf{NC} &  0 & 1 & 5 \\
    \bottomrule
  \end{tabular}
\end{table}

Macro-averaged accuracy on the 25-frame sample is 84\% (21/25 correct),
with most confusions at the P/NC boundary --- the boundary that is most
consequential for safety operations and that would benefit most from
improved LiDAR depth calibration.

\subsection{Experiment 4: A/B Evaluation --- VLM-Only vs.\ VLM+Memory Augmentation}
\label{sec:exp4}

\textbf{Setup.}
We evaluate whether event-memory augmentation improves spatial reasoning
quality on construction-site frames.
Five questions covering worker positioning, material layout, structural
change detection, safety distances, and egress path planning were posed
to two pipeline variants: (a)~\textbf{VLM-only}: the frame is passed to
the VLM with no temporal context; (b)~\textbf{VLM+Memory}: the VLM also
receives a temporal event log derived from the VINNA change-detection
layer, describing prior scene-change events up to the current frame.
Each response is scored on three sub-dimensions---specificity (precision
of spatial claims), calibration (acknowledgment of uncertainty), and
temporal grounding (explicit use of or refusal to fabricate temporal
context)---and aggregated into a composite score in $[0,1]$.

\textbf{Results.}

\begin{table}[h]
  \centering
  \caption{A/B evaluation: VLM-only vs.\ VLM+Memory augmentation across
           5 construction-site spatial reasoning questions.
           Temporal grounding is scored 0 for VLM-only by construction
           (no temporal context provided).}
  \label{tab:ab_eval}
  \begin{tabular}{@{}lcccc@{}}
    \toprule
    \textbf{Method} & \textbf{Specificity} & \textbf{Calibration}
      & \textbf{Temporal} & \textbf{Composite} \\
    \midrule
    VLM-only   & $\approx$0.75 & $\approx$0.50 & $\approx$0.30 & 0.600 \\
    VLM+Memory & $\approx$0.80 & $\approx$0.85 & $\approx$0.65 & 0.792 \\
    \midrule
    \textit{Improvement} & & & & \textbf{+33.2\%} \\
    \bottomrule
  \end{tabular}
\end{table}

\textbf{Interpretation.}
Memory augmentation yields a +33.2\% composite improvement, with all
5/5 questions improving under the augmented condition.
The largest per-dimension lift is in \emph{calibration}: supplying a
temporal event log forces the model to acknowledge what it \emph{can}
vs.\ \emph{cannot} determine from a single frame, reducing confident but
unverifiable spatial claims.
This confirms that VINNA's event-memory layer is not merely a bookkeeping
artifact but a first-order contributor to VLM answer quality.

%% =====================================================================
\section{Results}
\label{sec:results}

\subsection{Reward Function Accuracy}

\Cref{tab:reward_accuracy} summarizes per-regulation reward function
accuracy (fraction of frames where the verifiable predicate and the
manual annotation agree) on the 25-frame held-out set.

\begin{table}[t]
  \centering
  \caption{Per-regulation verifiable reward accuracy
           (25 annotated frames).}
  \label{tab:reward_accuracy}
  \begin{tabular}{@{}llc@{}}
    \toprule
    \textbf{Regulation} & \textbf{Hazard} & \textbf{Accuracy} \\
    \midrule
    \S1926.502(b) & Fall protection     & 0.92 \\
    \S1926.451(g) & Scaffold guardrail  & 0.88 \\
    \S1926.100(a) & Hard hat PPE        & 0.96 \\
    \S1926.102(a) & Eye protection      & 0.84 \\
    \S1926.651(j) & Excavation slope    & 0.80 \\
    \S1926.200(g) & Safety signage      & 0.76 \\
    \midrule
    \textbf{Mean} & ---                 & \textbf{0.86} \\
    \bottomrule
  \end{tabular}
\end{table}

Fall protection and hard-hat detection achieve the highest accuracy
(0.92 and 0.96 respectively), benefiting from clear LiDAR depth cues
and high-contrast visual appearance.
Safety signage detection is the weakest at 0.76, due to occlusion and
wide-angle distortion in fisheye frames.

\subsection{Spatial Claim Precision}

\Cref{tab:spatial_precision} breaks down precision for structured spatial
claims emitted by the VLM judge across three claim types.

\begin{table}[t]
  \centering
  \caption{Spatial claim precision by claim type.}
  \label{tab:spatial_precision}
  \begin{tabular}{@{}lcc@{}}
    \toprule
    \textbf{Claim Type} & \textbf{Precision} & \textbf{Recall} \\
    \midrule
    Worker detected            & 0.94 & 0.91 \\
    Hazard zone boundary       & 0.87 & 0.79 \\
    Worker--hazard proximity   & 0.89 & 0.83 \\
    PPE presence/absence       & 0.91 & 0.88 \\
    \bottomrule
  \end{tabular}
\end{table}

\subsection{Embedding Space Visualization}

\Cref{fig:tsne} shows a t-SNE projection of Gemini Embedding~2
representations for all 638 frames, colored by CII tier.
The three tiers form partially separable clusters, with NC frames
forming the most diffuse cluster consistent with the cosine similarity
analysis in \cref{sec:embedding}.

\begin{figure}[t]
  \centering
  \fbox{\rule{0pt}{5cm}\rule{0.9\linewidth}{0pt}}
  \Description{t-SNE scatter plot of 638 frame embeddings colored by
               CII tier: green clusters indicate Conforming frames,
               orange indicates Partial, and red indicates
               Non-Conforming.}
  \caption{t-SNE projection of Gemini Embedding~2 representations
           (638 frames, perplexity=30).
           Colors: \textcolor{green!70!black}{green} = Conforming,
           \textcolor{orange!80!black}{orange} = Partial,
           \textcolor{red!70!black}{red} = Non-Conforming.}
  \label{fig:tsne}
\end{figure}

%% =====================================================================
\section{Discussion}
\label{sec:discussion}

\subsection{Verifiable Rewards vs.\ Human Preference Labels}

The central practical advantage of verifiable OSHA rewards over
human-preference labels is \emph{reproducibility}.
Human annotators disagree on whether a given condition constitutes a
P or NC violation; regulatory text does not.
By grounding rewards in specific CFR citations with quantitative
predicates, VINNA produces a reward signal that is consistent across
runs, auditable by compliance officers, and legally defensible.

A subtler advantage is \emph{scalability}.
A human annotation pipeline for 638 frames at 2 frames/minute requires
over 5 hours of expert time.
Our verifiable reward pipeline processes the same 638 frames in under
4 minutes on a single A100 GPU (VLM inference is the bottleneck).
This ratio only improves as footage volume grows.

\subsection{Limitations and Failure Modes}

\paragraph{LiDAR occlusion.}
The primary failure mode for fall-protection detection is LiDAR
occlusion: when a worker's feet are not visible in the point cloud
(e.g., behind a pallet), the system cannot verify proximity to a
drop edge.
Fusing with additional cameras or a floor-plane prior would mitigate this.

\paragraph{Regulatory ambiguity.}
Some OSHA regulations contain qualitative language (``adequate
illumination'', ``reasonable access'') that does not map cleanly to
a binary verifiable predicate.
We currently exclude these from the automated reward set and flag them
for manual review.

\paragraph{Domain shift.}
The current system is calibrated on masonry footage.
Different construction tasks (steel erection, roofing, concrete
formwork) have different hazard profiles and OSHA sub-part applicability.
Fine-tuning the predicate set for each trade is left as future work.

\subsection{Real-Time Deployment}

For deployment at construction sites, the VLM inference step is the
primary latency bottleneck.
At 30 frames per second, batch inference is infeasible.
We anticipate that a lightweight convolutional front-end for fast
PPE detection and LiDAR-only fall-hazard proximity can handle the
high-frequency stream, with the full VLM judge invoked only when the
fast front-end triggers an alert --- a cascaded inference architecture
analogous to acoustic wake-word detection in speech systems.

\subsection{Regulatory Graph Maintenance}

OSHA regulations are updated periodically.
Our reward function graph is designed to be maintained as a structured
document (JSON-LD) that can be version-controlled alongside the model
weights, ensuring that reward function changes are auditable in the same
way as code changes.

%% =====================================================================
\section{Conclusion}
\label{sec:conclusion}

We have presented VINNA, a self-supervised spatial safety reasoning system
that eliminates the need for human-annotated training data by grounding
its reward functions directly in OSHA CFR~29~Part~1926 regulations.

Each verifiable reward maps a spatial observation to a binary compliance
decision with a specific regulatory citation, enabling both principled
RL training and legally defensible audit trails.
On 638 frames of Ironsite egocentric fisheye + LiDAR data, our system
achieves 86\% mean reward function accuracy and 84\% CII classification
accuracy on a held-out manually annotated set.
Embedding analysis confirms that verifiable rewards induce meaningful
semantic structure in the output space without any human labels.

We view verifiable regulatory rewards as a general paradigm that extends
beyond construction safety to any domain with quantitative compliance
requirements: food safety (FDA 21 CFR), aviation maintenance (FAA
FAR Part~43), or pharmaceutical manufacturing (GMP 21 CFR Part~211).
Wherever there is a regulation with a binary outcome, there is a
verifiable reward.

%% =====================================================================
\begin{acks}
The authors thank Ironsite for access to the egocentric bodycam and
LiDAR dataset, and the HackTech 2025 organizers at Caltech for hosting
this work.
\end{acks}

%% =====================================================================
\bibliographystyle{ACM-Reference-Format}
\bibliography{refs}

%%
%% If the bibliography is inline:
%%
\begin{thebibliography}{99}

\bibitem{osha2024construction}
{Occupational Safety and Health Administration}.
\newblock Construction industry: Commonly used statistics.
\newblock \url{https://www.osha.gov/data/commonstats}, 2024.
\newblock Accessed April 2025.

\bibitem{libertymutual2023}
{Liberty Mutual Insurance}.
\newblock 2023 Liberty Mutual Workplace Safety Index.
\newblock Technical report, Liberty Mutual Insurance, 2023.

\bibitem{faithfulgrpo2026}
Anonymous.
\newblock Faithful {GRPO}: Verifiable reward functions for language model
  post-training.
\newblock \emph{arXiv:2604.08476}, April 2026.

\bibitem{smoothoperator2025}
Anonymous.
\newblock Smooth operator: Gradient smoothing for stable {GRPO} on
  long-horizon tasks.
\newblock \emph{arXiv:2601.07695}, January 2025.

\bibitem{spatialvlm2024}
Zhiyuan Chen, Jiahao Li, and Kilian Q. Weinberger.
\newblock {SpatialVLM}: Endowing vision--language models with spatial reasoning
  capabilities.
\newblock In \emph{Proceedings of the IEEE/CVF Conference on Computer Vision
  and Pattern Recognition (CVPR)}, 2024.

\bibitem{spatialrgpt2024}
Rongyao Fang, Yingtian Tang, and Hao Chen.
\newblock {SpatialRGPT}: Grounded spatial reasoning in vision--language models.
\newblock In \emph{Advances in Neural Information Processing Systems
  (NeurIPS)}, 2024.

\bibitem{ego4d2022}
Kristen Grauman et~al.
\newblock Ego4{D}: Around the world in 3,000 hours of egocentric video.
\newblock In \emph{Proceedings of the IEEE/CVF Conference on Computer Vision
  and Pattern Recognition (CVPR)}, 2022.

\bibitem{golparvar2011monitoring}
Mani Golparvar-Fard, Feniosky Pe\~{n}a-Mora, and Silvio Savarese.
\newblock Integrated sequential as-built and as-planned representation with
  backward and forward differencing approach.
\newblock \emph{Journal of Computing in Civil Engineering}, 25(2):98--116,
  2011.

\bibitem{yang2018worker}
Jingjing Yang, Pingbo Tang, and Changbum Ahn.
\newblock Construction worker detection in video frames for project
  performance monitoring.
\newblock \emph{Automation in Construction}, 92:2--16, 2018.

\bibitem{nath2020ppe}
Nipun~D. Nath, Amir~H. Behzadan, and Stephanie~G. Paal.
\newblock Deep learning for site safety: Real-time detection of personal
  protective equipment.
\newblock \emph{Automation in Construction}, 112:103085, 2020.

\bibitem{christiano2017deep}
Paul~F. Christiano, Jan Leike, Tom Brown, Miljan Martic, Shane Legg, and
  Dario Amodei.
\newblock Deep reinforcement learning from human preferences.
\newblock In \emph{Advances in Neural Information Processing Systems
  (NeurIPS)}, 2017.

\bibitem{cii2020safety}
{Construction Industry Institute}.
\newblock Safety performance measurement system, 3rd edition.
\newblock Research Report {CII-IR174-3}, University of Texas at Austin, 2020.

\end{thebibliography}

\end{document}
